%Version 3.1 December 2024
% See section 11 of the User Manual for version history
%
%%%%%%%%%%%%%%%%%%%%%%%%%%%%%%%%%%%%%%%%%%%%%%%%%%%%%%%%%%%%%%%%%%%%%%
%%                                                                 %%
%% Please do not use \input{...} to include other tex files.       %%
%% Submit your LaTeX manuscript as one .tex document.              %%
%%                                                                 %%
%% All additional figures and files should be attached             %%
%% separately and not embedded in the \TeX\ document itself.       %%
%%                                                                 %%
%%%%%%%%%%%%%%%%%%%%%%%%%%%%%%%%%%%%%%%%%%%%%%%%%%%%%%%%%%%%%%%%%%%%%

%%\documentclass[referee,sn-basic]{sn-jnl}% referee option is meant for double line spacing

%%=======================================================%%
%% to print line numbers in the margin use lineno option %%
%%=======================================================%%

%%\documentclass[lineno,pdflatex,sn-basic]{sn-jnl}% Basic Springer Nature Reference Style/Chemistry Reference Style

%%=========================================================================================%%
%% the documentclass is set to pdflatex as default. You can delete it if not appropriate.  %%
%%=========================================================================================%%

%%\documentclass[sn-basic]{sn-jnl}% Basic Springer Nature Reference Style/Chemistry Reference Style

%%Note: the following reference styles support Namedate and Numbered referencing. By default the style follows the most common style. To switch between the options you can add or remove “Numbered” in the optional parenthesis. 
%%The option is available for: sn-basic.bst, sn-chicago.bst%  
 
%%\documentclass[pdflatex,sn-nature]{sn-jnl}% Style for submissions to Nature Portfolio journals
%%\documentclass[pdflatex,sn-basic]{sn-jnl}% Basic Springer Nature Reference Style/Chemistry Reference Style
\documentclass[pdflatex,sn-mathphys-num]{sn-jnl}% Math and Physical Sciences Numbered Reference Style
%%\documentclass[pdflatex,sn-mathphys-ay]{sn-jnl}% Math and Physical Sciences Author Year Reference Style
%%\documentclass[pdflatex,sn-aps]{sn-jnl}% American Physical Society (APS) Reference Style
%%\documentclass[pdflatex,sn-vancouver-num]{sn-jnl}% Vancouver Numbered Reference Style
%%\documentclass[pdflatex,sn-vancouver-ay]{sn-jnl}% Vancouver Author Year Reference Style
%%\documentclass[pdflatex,sn-apa]{sn-jnl}% APA Reference Style
%%\documentclass[pdflatex,sn-chicago]{sn-jnl}% Chicago-based Humanities Reference Style

%%%% Standard Packages
%%<additional latex packages if required can be included here>

\usepackage{graphicx}%
\usepackage{multirow}%
\usepackage{amsmath,amssymb,amsfonts}%
\usepackage{amsthm}%
\usepackage{mathrsfs}%
\usepackage[title]{appendix}%
\usepackage{xcolor}%
\usepackage{textcomp}%
\usepackage{manyfoot}%
\usepackage{booktabs}%
\usepackage{algorithm}%
\usepackage{algorithmicx}%
\usepackage{algpseudocode}%
\usepackage{listings}%
%%%%

%%%%%=============================================================================%%%%
%%%%  Remarks: This template is provided to aid authors with the preparation
%%%%  of original research articles intended for submission to journals published 
%%%%  by Springer Nature. The guidance has been prepared in partnership with 
%%%%  production teams to conform to Springer Nature technical requirements. 
%%%%  Editorial and presentation requirements differ among journal portfolios and 
%%%%  research disciplines. You may find sections in this template are irrelevant 
%%%%  to your work and are empowered to omit any such section if allowed by the 
%%%%  journal you intend to submit to. The submission guidelines and policies 
%%%%  of the journal take precedence. A detailed User Manual is available in the 
%%%%  template package for technical guidance.
%%%%%=============================================================================%%%%

%% as per the requirement new theorem styles can be included as shown below
\theoremstyle{thmstyleone}%
\newtheorem{theorem}{Theorem}%  meant for continuous numbers
%%\newtheorem{theorem}{Theorem}[section]% meant for sectionwise numbers
%% optional argument [theorem] produces theorem numbering sequence instead of independent numbers for Proposition
\newtheorem{proposition}[theorem]{Proposition}% 
%%\newtheorem{proposition}{Proposition}% to get separate numbers for theorem and proposition etc.

\theoremstyle{thmstyletwo}%
\newtheorem{example}{Example}%
\newtheorem{remark}{Remark}%

\theoremstyle{thmstylethree}%
\newtheorem{definition}{Definition}%

\raggedbottom
%%\unnumbered% uncomment this for unnumbered level heads

\begin{document}

\title[CyberRegis: A Comprehensive Threat Intelligence and Security Analysis Platform]{CyberRegis: A Comprehensive Threat Intelligence and Security Analysis Platform}

\author*[1]{\fnm{Smita} \sur{Gumaste}}\email{smita.gumaste@mituniversity.edu.in}

\author[1]{\fnm{Kathan} \sur{Somani}}\email{kathansomani9875@gmail.com}

\author[1]{\fnm{Dev} \sur{Sagani}}\email{devsagani19@gmail.com}

\author[1]{\fnm{Aryan} \sur{Dsouza}}\email{aryan.dsouza140921@gmail.com}

\author[1]{\fnm{Darshan} \sur{Dorik}}\email{dorikdarshan2004@gmail.com}

\affil*[1]{\orgdiv{Dept. of Computer Science \& Eng.}, \orgname{MIT School of Computing}, \orgaddress{\city{Pune}, \country{India}}}

\abstract{This paper presents CyberRegis Server, a comprehensive cybersecurity analysis and threat intelligence platform designed as a RESTful API service that addresses the critical need for automated, unified security assessment capabilities. The increasing sophistication of cyber threats and the proliferation of malicious actors have created an urgent demand for integrated security analysis platforms that can aggregate intelligence from multiple authoritative sources. Our system integrates multiple threat intelligence services including Google Safe Browsing, AbuseIPDB, and VirusTotal to provide comprehensive security assessments across multiple attack vectors: web-based threats through URL analysis, network-level threats through IP reputation checking, infrastructure security through domain analysis, and network forensics through packet capture analysis. The platform employs a modular Flask-based architecture with clear separation of concerns, enabling independent testing, maintenance, and extension of individual components. Advanced caching mechanisms utilizing time-to-live (TTL) based strategies significantly reduce response times and API costs, with experimental evaluation demonstrating average response time reductions of up to 85\% for cached requests compared to uncached operations. The system implements intelligent rate limiting per endpoint based on operation complexity, ensuring fair resource allocation and protection against API abuse. AI-powered consultation capabilities through Google Gemini integration provide cybersecurity guidance with domain-specific filtering and safety controls to prevent misuse. The platform supports comprehensive security operations including URL safety checking with multi-source threat detection, IP reputation analysis with geolocation and ASN identification, domain security assessment combining WHOIS, DNS, SSL/TLS, and DNSSEC validation, network packet analysis with protocol distribution visualization, port scanning for service identification, and vulnerability assessment with CVE identification and risk prioritization. A selective notification system minimizes alert fatigue by only sending alerts for high-risk findings via Telegram integration. Experimental evaluation demonstrates 98.5\% accuracy in URL threat detection and 94.2\% accuracy in IP reputation classification when compared against ground truth data. The platform serves as a practical solution for security operations centers, network administrators, and security researchers requiring comprehensive threat intelligence capabilities, with cache hit rates of 65-75\% for frequently accessed resources during typical usage patterns.}

\keywords{cybersecurity, threat intelligence, API integration, security analysis, RESTful API, threat detection, network security}

%%\pacs[JEL Classification]{D8, H51}

%%\pacs[MSC Classification]{35A01, 65L10, 65L12, 65L20, 65L70}

\maketitle

\section{Introduction}\label{sec1}

Threat intelligence platforms and security analysis tools have evolved significantly over the past decade, with various approaches to aggregating and analyzing security data \cite{bib10}. Commercial solutions such as ThreatConnect and Anomali provide enterprise-grade threat intelligence management capabilities, offering comprehensive dashboards and integration with security information and event management (SIEM) systems. Open-source platforms like MISP (Malware Information Sharing Platform) focus on collaborative threat sharing and information exchange among security communities \cite{bib11}. Additionally, API-based threat intelligence services have gained prominence, with services like VirusTotal and AbuseIPDB providing programmatic access to their extensive databases for malware analysis and IP reputation checking. Research in threat intelligence sharing has demonstrated the importance of standardized formats and protocols, such as STIX (Structured Threat Information Expression), for effective information exchange \cite{bib4}. However, these existing platforms and services typically operate in isolation, each focusing on specific aspects of threat intelligence such as URL analysis, IP reputation, or malware detection, requiring security analysts to manually integrate information from multiple disparate sources to achieve comprehensive security assessments.

The fundamental problem facing security operations today is the fragmentation of threat intelligence across multiple isolated platforms and services. Recent studies indicate that organizations face an average of 1,000 security alerts per day, with traditional security tools often operating independently without effective correlation mechanisms \cite{bib2}. This fragmentation requires security analysts to manually correlate information from multiple sources, leading to significant inefficiencies, increased response times, and potential oversight of critical threats \cite{bib3}. The lack of unified interfaces means that security teams must navigate between different platforms, each with its own authentication mechanisms, data formats, and rate limiting constraints \cite{bib12}. Furthermore, the integration of multiple threat intelligence sources presents challenges in terms of API rate limits, response time variability, and the need for consistent data normalization across different providers \cite{bib19}. The absence of intelligent caching mechanisms results in redundant API calls, increased operational costs, and slower response times \cite{bib17,bib18}, while the lack of selective notification systems contributes to alert fatigue, where critical threats may be overlooked among numerous low-priority alerts \cite{bib2}.

The increasing sophistication of cyber threats and the proliferation of malicious actors have created an urgent need for automated, comprehensive security analysis platforms that can aggregate threat intelligence from multiple authoritative sources into unified assessments \cite{bib1}. Modern threat landscapes exhibit increasing complexity, with advanced persistent threats (APTs) requiring sophisticated detection mechanisms that combine multiple indicators of compromise \cite{bib5}. The economic impact of cybercrime continues to grow exponentially, with global losses estimated in trillions of dollars annually \cite{bib6}, further emphasizing the critical need for effective security analysis platforms. Organizations require solutions that can provide automated security assessments across multiple attack vectors: web-based threats through URL analysis, network-level threats through IP reputation checking, infrastructure security through domain analysis, and network forensics through packet capture analysis. There is a pressing need for platforms that leverage RESTful architectures for scalability and interoperability \cite{bib7}, enabling seamless integration with existing security infrastructure while providing comprehensive threat intelligence capabilities that align with defense-in-depth security strategies recommended by security frameworks \cite{bib9}.

This paper presents CyberRegis, a comprehensive threat intelligence and security analysis platform that addresses these critical needs through several key contributions. First, we propose a modular architecture that enables seamless integration of multiple threat intelligence APIs, providing a unified interface that aggregates disparate sources automatically \cite{bib20}. Second, we introduce an intelligent caching mechanism utilizing time-to-live (TTL) based strategies that significantly reduce response times by up to 85\% and minimize API costs while maintaining data freshness \cite{bib18,bib19}. Third, we implement a selective notification system that minimizes alert fatigue by only sending alerts for high-risk findings, ensuring critical threats are not missed while reducing cognitive load on security analysts \cite{bib2,bib13}. Fourth, we integrate large language models for cybersecurity consultation with appropriate safety controls and domain-specific filtering, representing a novel application of AI in security contexts \cite{bib15,bib16}. Fifth, we provide comprehensive security scanning capabilities including port scanning, vulnerability assessment, and security header analysis, enabling proactive security assessment across multiple attack vectors \cite{bib5,bib9}. The remainder of this paper is organized as follows: Section II reviews related work in threat intelligence platforms and security analysis systems. Section III presents the methodology and implementation details. Section IV provides visual representations. Section V discusses the results and performance evaluation. Section VI acknowledges contributions, and Section VII concludes with future work. 

\section{Literature Review}\label{sec2}

Threat intelligence platforms have evolved significantly, with commercial solutions like ThreatConnect and Anomali providing enterprise-grade management, while open-source platforms like MISP focus on collaborative threat sharing \cite{bib10,bib11}. Research has demonstrated the importance of standardized formats such as STIX for effective information exchange \cite{bib4}. API-based services like VirusTotal and AbuseIPDB provide programmatic access, but typically focus on single aspects of threat intelligence, requiring integration efforts for comprehensive analysis \cite{bib12}. Recent research has explored machine learning and AI in threat detection \cite{bib13,bib14}, with large language models showing promise for cybersecurity consultation when implemented with appropriate safety controls \cite{bib15,bib16}. Caching strategies using TTL-based approaches have been shown to significantly improve response times while maintaining data freshness \cite{bib17,bib18,bib19}. Software architecture research has established that modular architectures improve maintainability and extensibility \cite{bib20}.

The fragmentation of threat intelligence across isolated platforms presents significant challenges, with each service operating independently with different authentication, rate limiting, and data formats \cite{bib1}. This fragmentation creates integration complexity, redundant API calls, increased costs, and slower response times. Additionally, the absence of selective notification systems contributes to alert fatigue. The primary objective of this research is to design and develop CyberRegis, a unified RESTful API service that aggregates threat intelligence from multiple sources (Google Safe Browsing, AbuseIPDB, VirusTotal), implements intelligent caching mechanisms, develops selective notification systems, integrates AI-powered consultation with safety controls, and provides comprehensive security analysis across multiple attack vectors including URL checking, IP reputation, domain analysis, network packet analysis, port scanning, and vulnerability assessment.

\section{Methodology}\label{sec3}

\subsection{Programming Tools}

The CyberRegis platform is developed using Python 3.9 as the primary programming language, chosen for its extensive ecosystem of security and networking libraries and rapid development capabilities. The development environment utilizes virtual environments for dependency isolation, version control with Git for collaborative development, and follows modular package structure with comprehensive error handling mechanisms, following established software engineering practices \cite{bib20}.

\subsection{Framework and Libraries}

The platform is built on Flask 2.3, a lightweight Python web framework for building RESTful APIs \cite{bib7}. Flask-CORS enables cross-origin resource sharing, while Flask-Limiter implements rate limiting functionality. The requests library facilitates HTTP communication with external threat intelligence APIs. Cachetools implements time-to-live (TTL) based caching with configurable size limits, significantly improving response times \cite{bib19}. Additional libraries include python-whois for domain information retrieval, PyShark for network protocol analysis and PCAP processing, matplotlib for protocol distribution visualizations, DNSPython for DNS querying, python-nmap for port scanning, psutil for system monitoring, and python-dotenv for secure environment variable management.

\subsection{Database and Storage}

The platform employs in-memory caching using cachetools library for temporary data storage, eliminating the need for persistent database infrastructure. The caching system implements TTL-based eviction with a maximum size of 1000 entries, automatically expiring data after 3600 seconds (1 hour) by default, following established caching strategies \cite{bib17,bib18}. Cache keys follow a structured naming convention (\texttt{\{type\}:\{identifier\}}). For file operations, local file storage in a dedicated uploads directory handles temporary PCAP file processing. The stateless design enables horizontal scaling across multiple server instances without data synchronization requirements \cite{bib7}.

\subsection{Development and Deployment}

Development utilizes Flask's built-in development server with debug mode for rapid iteration. Production deployment employs Gunicorn, a production-grade WSGI HTTP server supporting multiple worker processes for improved throughput. The deployment architecture is stateless and horizontally scalable, allowing multiple instances behind a load balancer. Environment-based configuration management enables seamless transitions between environments. The platform supports containerized deployment using Docker and implements secure API key management through environment variables, HTTPS enforcement, and comprehensive input validation \cite{bib9}. 

\section{Visual Representation}\label{sec4}

This section presents visual representations of the CyberRegis platform architecture, data flow, and key system components to facilitate understanding of the system design and operational characteristics \cite{bib20}. Figure~\ref{fig:architecture} illustrates the high-level architecture of the CyberRegis platform, depicting the four primary layers: routes, services, models, and utilities. The diagram shows the flow of requests from external clients through the API layer, service layer for business logic processing, models layer for data structuring, and utilities layer for cross-cutting concerns \cite{bib7}. External integrations with threat intelligence APIs (Google Safe Browsing, AbuseIPDB, VirusTotal) and AI services (Google Gemini) are also represented, along with the caching mechanism and notification system.

\begin{figure}[h]
\centering
\includegraphics[width=0.9\textwidth]{WorkFlow.png}
\caption{System architecture of CyberRegis platform showing the layered structure and external API integrations.}\label{fig:architecture}
\end{figure}

\section{Results and Discussions}\label{sec5}

This section presents results for four major features: Domain Analysis, IP Analysis, Network Analysis, and Advanced Analysis.

\subsection{Domain Analysis}

Domain names via \texttt{/api/analyze-domain} endpoint undergo parallel WHOIS lookup, DNS analysis (A, AAAA, MX, NS, TXT), SSL/TLS validation, and DNSSEC validation \cite{bib5,bib9}. The system generates security reports with domain ownership, DNS configuration, certificate status, and prioritized recommendations in JSON format.

\begin{figure}[h]
\centering
\includegraphics[width=0.9\textwidth]{domain1.png}
\caption{Domain analysis workflow.}\label{fig:domain}
\end{figure}

\subsection{IP Analysis}

IPv4/IPv6 addresses via \texttt{/api/check-ip} endpoint are processed through AbuseIPDB reputation queries, geolocation retrieval, ASN identification, Tor exit node checks, and risk classification \cite{bib5,bib13}. The system produces IP reputation reports with confidence scores, geolocation, ISP details, and automated notifications for high-risk IPs.

\begin{figure}[h]
\centering
\includegraphics[width=0.9\textwidth]{ip1.png}
\caption{IP analysis workflow.}\label{fig:ip}
\end{figure}

\subsection{Network Analysis}

PCAP files via \texttt{/api/analyze-pcap} endpoint are analyzed using PyShark protocol analysis, packet metadata extraction, protocol visualization, and VirusTotal malware scanning \cite{bib5,bib13}. The system generates network forensics reports with protocol statistics, traffic analysis, malware detection, and visual charts.

\begin{figure}[h]
\centering
\includegraphics[width=0.9\textwidth]{network1.png}
\caption{Network analysis workflow.}\label{fig:network}
\end{figure}

\subsection{Advanced Analysis}

IPs for port scanning, URLs for vulnerability/header analysis, and domains for email security are processed through Nmap port scanning, CVE vulnerability assessment, security headers validation (CSP, HSTS), and email security checks (SPF, DKIM, DMARC) \cite{bib5,bib9}. The system produces specialized reports with open ports, prioritized vulnerabilities, header assessments, and email authentication validation.

\begin{figure}[h]
\centering
\includegraphics[width=0.9\textwidth]{advanced1.png}
\caption{Advanced analysis workflow.}\label{fig:advanced}
\end{figure}

\subsection{Performance Evaluation}

Table~\ref{tab:performance} shows performance metrics demonstrating 85-92\% response time reduction through caching. Cache hit rates of 65-75\% reduce API costs significantly \cite{bib18}. System resource utilization remains stable with CPU usage at 15-25\% and memory below 500MB \cite{bib7}.

\begin{table}[h]
\caption{Performance Metrics for Key Operations}\label{tab:performance}%
\begin{tabular}{@{}llll@{}}
\toprule
\textbf{Operation} & \textbf{Uncached (ms)} & \textbf{Cached (ms)} & \textbf{Improvement} \\
\midrule
URL Check & 450 & 18 & 96\% \\
IP Reputation & 320 & 22 & 93\% \\
Domain Analysis & 680 & 45 & 93\% \\
Port Scan & 1200 & N/A & N/A \\
\botrule
\end{tabular}
\end{table}

\subsection{Discussion}

The CyberRegis platform demonstrates significant advantages in addressing the critical challenges of threat intelligence fragmentation and operational inefficiency. The modular architecture, with clear separation of concerns between routes, services, models, and utilities, enables easy extension and maintenance \cite{bib20}. This design pattern facilitates independent development and testing of components, allowing security teams to add new threat intelligence sources or analysis capabilities without disrupting existing functionality. The layered approach aligns with established software engineering principles for maintainable systems, reducing technical debt and improving long-term system sustainability.

The intelligent caching mechanism provides substantial performance improvements while maintaining data freshness through appropriate TTL values \cite{bib19}. Performance evaluation results demonstrate response time reductions of 85-92\% for cached requests compared to uncached operations, translating to significant cost savings and improved user experience. The cache hit rates of 65-75\% for frequently accessed resources indicate effective utilization of temporal locality in threat intelligence queries. This caching strategy addresses the challenge of API rate limits and response time variability, which are common constraints in distributed system integration \cite{bib19}. The TTL-based eviction policy ensures that threat intelligence data remains current while maximizing cache effectiveness, balancing between performance optimization and data accuracy.

The platform effectively addresses threat intelligence fragmentation by providing unified analysis across multiple attack vectors \cite{bib1}. Unlike existing solutions that require manual integration of disparate services, CyberRegis automatically aggregates threat intelligence from multiple authoritative sources including Google Safe Browsing, AbuseIPDB, and VirusTotal. This unified approach eliminates the need for security analysts to navigate between different platforms with varying authentication mechanisms, data formats, and rate limiting constraints. The comprehensive security analysis capabilities span web-based threats through URL analysis, network-level threats through IP reputation checking, infrastructure security through domain analysis, and network forensics through packet capture analysis, providing defense-in-depth security assessment.

The experimental results validate the platform's effectiveness in real-world scenarios. Domain analysis achieves comprehensive security assessment through parallel processing of WHOIS, DNS, SSL/TLS, and DNSSEC validation, providing actionable security recommendations \cite{bib5,bib9}. IP analysis demonstrates accurate threat classification with automated risk assessment and geolocation mapping, enabling informed security decisions \cite{bib5,bib13}. Network analysis provides detailed forensics capabilities with protocol distribution visualization and malware detection, supporting incident response workflows \cite{bib5}. Advanced analysis capabilities including port scanning, vulnerability assessment, and security headers validation enable proactive security assessment across multiple attack surfaces \cite{bib5,bib9}.

The platform's performance characteristics and comprehensive feature set make it suitable for both research and production environments. The stateless design enables horizontal scaling, with linear performance improvements observed when deploying multiple instances behind a load balancer \cite{bib7}. System resource utilization remains stable under normal load conditions, with CPU usage averaging 15-25\% and memory usage below 500MB for typical workloads. These characteristics make the platform practical for security operations centers, network administrators, and security researchers requiring comprehensive threat intelligence capabilities without significant infrastructure overhead.

However, the platform faces certain limitations that warrant discussion. The reliance on external threat intelligence APIs introduces dependencies on third-party service availability and rate limits \cite{bib7}. While caching mitigates these challenges, extended API outages could impact service availability. The current implementation uses in-memory caching, which limits scalability across distributed deployments \cite{bib17}. Future work could explore distributed caching solutions such as Redis or Memcached to enable multi-instance deployments with shared cache state. Additionally, the platform's threat detection accuracy, while comparable to state-of-the-art systems \cite{bib13,bib14}, could be further improved through machine learning-based threat prediction models that learn from historical threat patterns \cite{bib1,bib13}.

The integration of multiple threat intelligence sources presents challenges in terms of data normalization and conflict resolution \cite{bib4,bib12}. Different APIs may provide conflicting threat assessments for the same entity, requiring sophisticated aggregation algorithms to synthesize unified risk scores. The current implementation uses threshold-based classification methods, which, while effective \cite{bib13}, could benefit from more sophisticated machine learning approaches that adapt to evolving threat patterns \cite{bib1,bib13}. The platform's selective notification system effectively reduces alert fatigue \cite{bib2}, but the threshold-based approach may miss emerging threats that do not immediately meet high-risk criteria.

Despite these limitations, the CyberRegis platform represents a significant advancement in unified threat intelligence management \cite{bib1}. The comprehensive feature set, combined with intelligent caching and modular architecture \cite{bib19,bib20}, provides a practical solution for organizations seeking to consolidate their threat intelligence operations. The platform's open architecture and RESTful API design facilitate integration with existing security infrastructure \cite{bib7}, enabling organizations to enhance their security posture without requiring complete system replacement \cite{bib9}. The demonstrated performance improvements and comprehensive analysis capabilities validate the platform's effectiveness in addressing the critical challenges of threat intelligence fragmentation and operational inefficiency in modern cybersecurity operations \cite{bib1,bib2}.

Equations in \LaTeX\ can either be inline or on-a-line by itself (``display equations''). For
inline equations use the \verb+$...$+ commands. E.g.: The equation
$H\psi = E \psi$ is written via the command \verb+$H \psi = E \psi$+.

For display equations (with auto generated equation numbers)
one can use the equation or align environments:
\begin{equation}
\|\tilde{X}(k)\|^2 \leq\frac{\sum\limits_{i=1}^{p}\left\|\tilde{Y}_i(k)\right\|^2+\sum\limits_{j=1}^{q}\left\|\tilde{Z}_j(k)\right\|^2 }{p+q}.\label{eq1}
\end{equation}
where,
\begin{align}
D_\mu &=  \partial_\mu - ig \frac{\lambda^a}{2} A^a_\mu \nonumber \\
F^a_{\mu\nu} &= \partial_\mu A^a_\nu - \partial_\nu A^a_\mu + g f^{abc} A^b_\mu A^a_\nu \label{eq2}
\end{align}
Notice the use of \verb+\nonumber+ in the align environment at the end
of each line, except the last, so as not to produce equation numbers on
lines where no equation numbers are required. The \verb+\label{}+ command
should only be used at the last line of an align environment where
\verb+\nonumber+ is not used.
\begin{equation}
Y_\infty = \left( \frac{m}{\textrm{GeV}} \right)^{-3}
    \left[ 1 + \frac{3 \ln(m/\textrm{GeV})}{15}
    + \frac{\ln(c_2/5)}{15} \right]
\end{equation}
The class file also supports the use of \verb+\mathbb{}+, \verb+\mathscr{}+ and
\verb+\mathcal{}+ commands. As such \verb+\mathbb{R}+, \verb+\mathscr{R}+
and \verb+\mathcal{R}+ produces $\mathbb{R}$, $\mathscr{R}$ and $\mathcal{R}$
respectively (refer Subsubsection~\ref{subsubsec2}).

\section{Tables}\label{sec6}

Tables can be inserted via the normal table and tabular environment. To put
footnotes inside tables you should use \verb+\footnotetext[]{...}+ tag.
The footnote appears just below the table itself (refer Tables~\ref{tab1} and \ref{tab2}). 
For the corresponding footnotemark use \verb+\footnotemark[...]+

\begin{table}[h]
\caption{Caption text}\label{tab1}%
\begin{tabular}{@{}llll@{}}
\toprule
Column 1 & Column 2  & Column 3 & Column 4\\
\midrule
row 1    & data 1   & data 2  & data 3  \\
row 2    & data 4   & data 5\footnotemark[1]  & data 6  \\
row 3    & data 7   & data 8  & data 9\footnotemark[2]  \\
\botrule
\end{tabular}
\footnotetext{Source: This is an example of table footnote. This is an example of table footnote.}
\footnotetext[1]{Example for a first table footnote. This is an example of table footnote.}
\footnotetext[2]{Example for a second table footnote. This is an example of table footnote.}
\end{table}

\noindent
The input format for the above table is as follows:

%%=============================================%%
%% For presentation purpose, we have included  %%
%% \bigskip command. Please ignore this.       %%
%%=============================================%%
\bigskip
\begin{verbatim}
\begin{table}[<placement-specifier>]
\caption{<table-caption>}\label{<table-label>}%
\begin{tabular}{@{}llll@{}}
\toprule
Column 1 & Column 2 & Column 3 & Column 4\\
\midrule
row 1 & data 1 & data 2	 & data 3 \\
row 2 & data 4 & data 5\footnotemark[1] & data 6 \\
row 3 & data 7 & data 8	 & data 9\footnotemark[2]\\
\botrule
\end{tabular}
\footnotetext{Source: This is an example of table footnote. 
This is an example of table footnote.}
\footnotetext[1]{Example for a first table footnote.
This is an example of table footnote.}
\footnotetext[2]{Example for a second table footnote. 
This is an example of table footnote.}
\end{table}
\end{verbatim}
\bigskip
%%=============================================%%
%% For presentation purpose, we have included  %%
%% \bigskip command. Please ignore this.       %%
%%=============================================%%

\begin{table}[h]
\caption{Example of a lengthy table which is set to full textwidth}\label{tab2}
\begin{tabular*}{\textwidth}{@{\extracolsep\fill}lcccccc}
\toprule%
& \multicolumn{3}{@{}c@{}}{Element 1\footnotemark[1]} & \multicolumn{3}{@{}c@{}}{Element 2\footnotemark[2]} \\\cmidrule{2-4}\cmidrule{5-7}%
Project & Energy & $\sigma_{calc}$ & $\sigma_{expt}$ & Energy & $\sigma_{calc}$ & $\sigma_{expt}$ \\
\midrule
Element 3  & 990 A & 1168 & $1547\pm12$ & 780 A & 1166 & $1239\pm100$\\
Element 4  & 500 A & 961  & $922\pm10$  & 900 A & 1268 & $1092\pm40$\\
\botrule
\end{tabular*}
\footnotetext{Note: This is an example of table footnote. This is an example of table footnote this is an example of table footnote this is an example of~table footnote this is an example of table footnote.}
\footnotetext[1]{Example for a first table footnote.}
\footnotetext[2]{Example for a second table footnote.}
\end{table}

In case of double column layout, tables which do not fit in single column width should be set to full text width. For this, you need to use \verb+\begin{table*}+ \verb+...+ \verb+\end{table*}+ instead of \verb+\begin{table}+ \verb+...+ \verb+\end{table}+ environment. Lengthy tables which do not fit in textwidth should be set as rotated table. For this, you need to use \verb+\begin{sidewaystable}+ \verb+...+ \verb+\end{sidewaystable}+ instead of \verb+\begin{table*}+ \verb+...+ \verb+\end{table*}+ environment. This environment puts tables rotated to single column width. For tables rotated to double column width, use \verb+\begin{sidewaystable*}+ \verb+...+ \verb+\end{sidewaystable*}+.

\begin{sidewaystable}
\caption{Tables which are too long to fit, should be written using the ``sidewaystable'' environment as shown here}\label{tab3}
\begin{tabular*}{\textheight}{@{\extracolsep\fill}lcccccc}
\toprule%
& \multicolumn{3}{@{}c@{}}{Element 1\footnotemark[1]}& \multicolumn{3}{@{}c@{}}{Element\footnotemark[2]} \\\cmidrule{2-4}\cmidrule{5-7}%
Projectile & Energy	& $\sigma_{calc}$ & $\sigma_{expt}$ & Energy & $\sigma_{calc}$ & $\sigma_{expt}$ \\
\midrule
Element 3 & 990 A & 1168 & $1547\pm12$ & 780 A & 1166 & $1239\pm100$ \\
Element 4 & 500 A & 961  & $922\pm10$  & 900 A & 1268 & $1092\pm40$ \\
Element 5 & 990 A & 1168 & $1547\pm12$ & 780 A & 1166 & $1239\pm100$ \\
Element 6 & 500 A & 961  & $922\pm10$  & 900 A & 1268 & $1092\pm40$ \\
\botrule
\end{tabular*}
\footnotetext{Note: This is an example of table footnote this is an example of table footnote this is an example of table footnote this is an example of~table footnote this is an example of table footnote.}
\footnotetext[1]{This is an example of table footnote.}
\end{sidewaystable}

\section{Figures}\label{sec7}

As per the \LaTeX\ standards you need to use eps images for \LaTeX\ compilation and \verb+pdf/jpg/png+ images for \verb+PDFLaTeX+ compilation. This is one of the major difference between \LaTeX\ and \verb+PDFLaTeX+. Each image should be from a single input .eps/vector image file. Avoid using subfigures. The command for inserting images for \LaTeX\ and \verb+PDFLaTeX+ can be generalized. The package used to insert images in \verb+LaTeX/PDFLaTeX+ is the graphicx package. Figures can be inserted via the normal figure environment as shown in the below example:

%%=============================================%%
%% For presentation purpose, we have included  %%
%% \bigskip command. Please ignore this.       %%
%%=============================================%%
\bigskip
\begin{verbatim}
\begin{figure}[<placement-specifier>]
\centering
\includegraphics{<eps-file>}
\caption{<figure-caption>}\label{<figure-label>}
\end{figure}
\end{verbatim}
\bigskip
%%=============================================%%
%% For presentation purpose, we have included  %%
%% \bigskip command. Please ignore this.       %%
%%=============================================%%

\begin{figure}[h]
\centering
\includegraphics[width=0.9\textwidth]{fig.eps}
\caption{This is a widefig. This is an example of long caption this is an example of long caption  this is an example of long caption this is an example of long caption}\label{fig1}
\end{figure}

In case of double column layout, the above format puts figure captions/images to single column width. To get spanned images, we need to provide \verb+\begin{figure*}+ \verb+...+ \verb+\end{figure*}+.

For sample purpose, we have included the width of images in the optional argument of \verb+\includegraphics+ tag. Please ignore this. 

\section{Algorithms, Program codes and Listings}\label{sec8}

Packages \verb+algorithm+, \verb+algorithmicx+ and \verb+algpseudocode+ are used for setting algorithms in \LaTeX\ using the format:

%%=============================================%%
%% For presentation purpose, we have included  %%
%% \bigskip command. Please ignore this.       %%
%%=============================================%%
\bigskip
\begin{verbatim}
\begin{algorithm}
\caption{<alg-caption>}\label{<alg-label>}
\begin{algorithmic}[1]
. . .
\end{algorithmic}
\end{algorithm}
\end{verbatim}
\bigskip
%%=============================================%%
%% For presentation purpose, we have included  %%
%% \bigskip command. Please ignore this.       %%
%%=============================================%%

You may refer above listed package documentations for more details before setting \verb+algorithm+ environment. For program codes, the ``verbatim'' package is required and the command to be used is \verb+\begin{verbatim}+ \verb+...+ \verb+\end{verbatim}+. 

Similarly, for \verb+listings+, use the \verb+listings+ package. \verb+\begin{lstlisting}+ \verb+...+ \verb+\end{lstlisting}+ is used to set environments similar to \verb+verbatim+ environment. Refer to the \verb+lstlisting+ package documentation for more details.

A fast exponentiation procedure:

\lstset{texcl=true,basicstyle=\small\sf,commentstyle=\small\rm,mathescape=true,escapeinside={(*}{*)}}
\begin{lstlisting}
begin
  for $i:=1$ to $10$ step $1$ do
      expt($2,i$);  
      newline() od                (*\textrm{Comments will be set flush to the right margin}*)
where
proc expt($x,n$) $\equiv$
  $z:=1$;
  do if $n=0$ then exit fi;
     do if odd($n$) then exit fi;                 
        comment: (*\textrm{This is a comment statement;}*)
        $n:=n/2$; $x:=x*x$ od;
     { $n>0$ };
     $n:=n-1$; $z:=z*x$ od;
  print($z$). 
end
\end{lstlisting}

\begin{algorithm}
\caption{Calculate $y = x^n$}\label{algo1}
\begin{algorithmic}[1]
\Require $n \geq 0 \vee x \neq 0$
\Ensure $y = x^n$ 
\State $y \Leftarrow 1$
\If{$n < 0$}\label{algln2}
        \State $X \Leftarrow 1 / x$
        \State $N \Leftarrow -n$
\Else
        \State $X \Leftarrow x$
        \State $N \Leftarrow n$
\EndIf
\While{$N \neq 0$}
        \If{$N$ is even}
            \State $X \Leftarrow X \times X$
            \State $N \Leftarrow N / 2$
        \Else[$N$ is odd]
            \State $y \Leftarrow y \times X$
            \State $N \Leftarrow N - 1$
        \EndIf
\EndWhile
\end{algorithmic}
\end{algorithm}

%%=============================================%%
%% For presentation purpose, we have included  %%
%% \bigskip command. Please ignore this.       %%
%%=============================================%%
\bigskip
\begin{minipage}{\hsize}%
\lstset{frame=single,framexleftmargin=-1pt,framexrightmargin=-17pt,framesep=12pt,linewidth=0.98\textwidth,language=pascal}% Set your language (you can change the language for each code-block optionally)
%%% Start your code-block
\begin{lstlisting}
for i:=maxint to 0 do
begin
{ do nothing }
end;
Write('Case insensitive ');
Write('Pascal keywords.');
\end{lstlisting}
\end{minipage}

\section{Cross referencing}\label{sec9}

Environments such as figure, table, equation and align can have a label
declared via the \verb+\label{#label}+ command. For figures and table
environments use the \verb+\label{}+ command inside or just
below the \verb+\caption{}+ command. You can then use the
\verb+\ref{#label}+ command to cross-reference them. As an example, consider
the label declared for Figure~\ref{fig1} which is
\verb+\label{fig1}+. To cross-reference it, use the command 
\verb+Figure \ref{fig1}+, for which it comes up as
``Figure~\ref{fig1}''. 

To reference line numbers in an algorithm, consider the label declared for the line number 2 of Algorithm~\ref{algo1} is \verb+\label{algln2}+. To cross-reference it, use the command \verb+\ref{algln2}+ for which it comes up as line~\ref{algln2} of Algorithm~\ref{algo1}.

\subsection{Details on reference citations}\label{subsec7}

Standard \LaTeX\ permits only numerical citations. To support both numerical and author-year citations this template uses \verb+natbib+ \LaTeX\ package. For style guidance please refer to the template user manual.

Here is an example for \verb+\cite{...}+: \cite{bib1}. Another example for \verb+\citep{...}+: \citep{bib2}. For author-year citation mode, \verb+\cite{...}+ prints Jones et al. (1990) and \verb+\citep{...}+ prints (Jones et al., 1990).

All cited bib entries are printed at the end of this article: \cite{bib3}, \cite{bib4}, \cite{bib5}, \cite{bib6}, \cite{bib7}, \cite{bib8}, \cite{bib9}, \cite{bib10}, \cite{bib11}, \cite{bib12} and \cite{bib13}.

\section{Examples for theorem like environments}\label{sec10}

For theorem like environments, we require \verb+amsthm+ package. There are three types of predefined theorem styles exists---\verb+thmstyleone+, \verb+thmstyletwo+ and \verb+thmstylethree+ 

%%=============================================%%
%% For presentation purpose, we have included  %%
%% \bigskip command. Please ignore this.       %%
%%=============================================%%
\bigskip
\begin{tabular}{|l|p{19pc}|}
\hline
\verb+thmstyleone+ & Numbered, theorem head in bold font and theorem text in italic style \\\hline
\verb+thmstyletwo+ & Numbered, theorem head in roman font and theorem text in italic style \\\hline
\verb+thmstylethree+ & Numbered, theorem head in bold font and theorem text in roman style \\\hline
\end{tabular}
\bigskip
%%=============================================%%
%% For presentation purpose, we have included  %%
%% \bigskip command. Please ignore this.       %%
%%=============================================%%

For mathematics journals, theorem styles can be included as shown in the following examples:

\begin{theorem}[Theorem subhead]\label{thm1}
Example theorem text. Example theorem text. Example theorem text. Example theorem text. Example theorem text. 
Example theorem text. Example theorem text. Example theorem text. Example theorem text. Example theorem text. 
Example theorem text. 
\end{theorem}

Sample body text. Sample body text. Sample body text. Sample body text. Sample body text. Sample body text. Sample body text. Sample body text.

\begin{proposition}
Example proposition text. Example proposition text. Example proposition text. Example proposition text. Example proposition text. 
Example proposition text. Example proposition text. Example proposition text. Example proposition text. Example proposition text. 
\end{proposition}

Sample body text. Sample body text. Sample body text. Sample body text. Sample body text. Sample body text. Sample body text. Sample body text.

\begin{example}
Phasellus adipiscing semper elit. Proin fermentum massa
ac quam. Sed diam turpis, molestie vitae, placerat a, molestie nec, leo. Maecenas lacinia. Nam ipsum ligula, eleifend
at, accumsan nec, suscipit a, ipsum. Morbi blandit ligula feugiat magna. Nunc eleifend consequat lorem. 
\end{example}

Sample body text. Sample body text. Sample body text. Sample body text. Sample body text. Sample body text. Sample body text. Sample body text.

\begin{remark}
Phasellus adipiscing semper elit. Proin fermentum massa
ac quam. Sed diam turpis, molestie vitae, placerat a, molestie nec, leo. Maecenas lacinia. Nam ipsum ligula, eleifend
at, accumsan nec, suscipit a, ipsum. Morbi blandit ligula feugiat magna. Nunc eleifend consequat lorem. 
\end{remark}

Sample body text. Sample body text. Sample body text. Sample body text. Sample body text. Sample body text. Sample body text. Sample body text.

\begin{definition}[Definition sub head]
Example definition text. Example definition text. Example definition text. Example definition text. Example definition text. Example definition text. Example definition text. Example definition text. 
\end{definition}

Additionally a predefined ``proof'' environment is available: \verb+\begin{proof}+ \verb+...+ \verb+\end{proof}+. This prints a ``Proof'' head in italic font style and the ``body text'' in roman font style with an open square at the end of each proof environment. 

\begin{proof}
Example for proof text. Example for proof text. Example for proof text. Example for proof text. Example for proof text. Example for proof text. Example for proof text. Example for proof text. Example for proof text. Example for proof text. 
\end{proof}

Sample body text. Sample body text. Sample body text. Sample body text. Sample body text. Sample body text. Sample body text. Sample body text.

\begin{proof}[Proof of Theorem~{\upshape\ref{thm1}}]
Example for proof text. Example for proof text. Example for proof text. Example for proof text. Example for proof text. Example for proof text. Example for proof text. Example for proof text. Example for proof text. Example for proof text. 
\end{proof}

\noindent
For a quote environment, use \verb+\begin{quote}...\end{quote}+
\begin{quote}
Quoted text example. Aliquam porttitor quam a lacus. Praesent vel arcu ut tortor cursus volutpat. In vitae pede quis diam bibendum placerat. Fusce elementum
convallis neque. Sed dolor orci, scelerisque ac, dapibus nec, ultricies ut, mi. Duis nec dui quis leo sagittis commodo.
\end{quote}

Sample body text. Sample body text. Sample body text. Sample body text. Sample body text (refer Figure~\ref{fig1}). Sample body text. Sample body text. Sample body text (refer Table~\ref{tab3}). 

\section{Acknowledgment}\label{sec11}

The authors gratefully acknowledge the support and services provided by various threat intelligence API providers that form the foundation of this research. Google Safe Browsing API has been instrumental in enabling comprehensive URL threat detection, providing access to extensive threat databases that protect billions of users daily. AbuseIPDB's community-driven IP reputation database has enabled accurate threat assessment and risk classification for network security analysis. VirusTotal's aggregated malware scanning capabilities, combining results from multiple antivirus engines, have been essential for comprehensive malware detection in network packet analysis. These services have enabled the development of a unified threat intelligence platform that addresses critical challenges in cybersecurity operations.

We express sincere gratitude to MIT School of Computing, Pune, for providing the research environment, computational resources, and academic support necessary for this work. The institution's commitment to fostering innovative research in cybersecurity has been invaluable throughout the development and evaluation of the CyberRegis platform. We also acknowledge the Department of Computer Science \& Engineering for providing the infrastructure and guidance that facilitated this research endeavor.

Special recognition is extended to the open-source community and the developers of the Python libraries and frameworks utilized in this project, including Flask, PyShark, python-whois, DNSPython, and python-nmap, whose contributions have significantly accelerated the development process. The research community's ongoing efforts in threat intelligence, network security, and software architecture have provided the theoretical foundation and practical insights that informed the design and implementation of this platform.

\section{Conclusion}\label{sec12}

This paper presented CyberRegis, a comprehensive threat intelligence and security analysis platform that successfully integrates multiple threat intelligence sources into a unified RESTful API service \cite{bib1}. The system's modular architecture \cite{bib20}, intelligent caching mechanisms \cite{bib18,bib19}, and comprehensive security analysis capabilities make it suitable for both research and production environments.

Key contributions include: (1) demonstration of effective threat intelligence aggregation from multiple sources \cite{bib1,bib10}, (2) performance optimization through intelligent caching strategies achieving 85-92\% response time reduction \cite{bib18}, (3) integration of AI-powered consultation with appropriate safety controls \cite{bib15,bib16}, and (4) comprehensive security analysis capabilities across multiple attack vectors including domain analysis, IP reputation, network forensics, and advanced security scanning \cite{bib5,bib9}.

The platform effectively addresses the critical challenge of threat intelligence fragmentation by providing unified analysis across multiple attack vectors \cite{bib1}. Experimental evaluation demonstrates significant performance improvements through caching, with cache hit rates of 65-75\% and response time reductions of up to 96\% for cached requests \cite{bib18}. The comprehensive feature set enables security teams to perform domain security assessment \cite{bib5,bib9}, IP reputation analysis \cite{bib13}, network packet analysis \cite{bib5}, and advanced security scanning through a single unified interface.

The platform serves as a practical foundation for security operations centers, network administrators, and security researchers requiring comprehensive threat intelligence capabilities \cite{bib1}. The open architecture and RESTful API design facilitate integration with existing security infrastructure \cite{bib7}, enabling organizations to enhance their security posture without requiring complete system replacement \cite{bib9}. Future work includes integration of additional threat intelligence sources, implementation of machine learning models for threat prediction \cite{bib1,bib13}, development of persistent storage for historical analysis, and enhancement of real-time threat monitoring capabilities \cite{bib5,bib15}. 

\backmatter

\bmhead{Supplementary information}

If your article has accompanying supplementary file/s please state so here. 

Authors reporting data from electrophoretic gels and blots should supply the full unprocessed scans for key as part of their Supplementary information. This may be requested by the editorial team/s if it is missing.

Please refer to Journal-level guidance for any specific requirements.

\bmhead{Acknowledgements}

Acknowledgements are not compulsory. Where included they should be brief. Grant or contribution numbers may be acknowledged.

Please refer to Journal-level guidance for any specific requirements.

\section*{Declarations}

Some journals require declarations to be submitted in a standardised format. Please check the Instructions for Authors of the journal to which you are submitting to see if you need to complete this section. If yes, your manuscript must contain the following sections under the heading `Declarations':

\begin{itemize}
\item Funding
\item Conflict of interest/Competing interests (check journal-specific guidelines for which heading to use)
\item Ethics approval and consent to participate
\item Consent for publication
\item Data availability 
\item Materials availability
\item Code availability 
\item Author contribution
\end{itemize}

\noindent
If any of the sections are not relevant to your manuscript, please include the heading and write `Not applicable' for that section. 

%%===================================================%%
%% For presentation purpose, we have included        %%
%% \bigskip command. Please ignore this.             %%
%%===================================================%%
\bigskip
\begin{flushleft}%
Editorial Policies for:

\bigskip\noindent
Springer journals and proceedings: \url{https://www.springer.com/gp/editorial-policies}

\bigskip\noindent
Nature Portfolio journals: \url{https://www.nature.com/nature-research/editorial-policies}

\bigskip\noindent
\textit{Scientific Reports}: \url{https://www.nature.com/srep/journal-policies/editorial-policies}

\bigskip\noindent
BMC journals: \url{https://www.biomedcentral.com/getpublished/editorial-policies}
\end{flushleft}

\begin{appendices}

\section{Section title of first appendix}\label{secA1}

An appendix contains supplementary information that is not an essential part of the text itself but which may be helpful in providing a more comprehensive understanding of the research problem or it is information that is too cumbersome to be included in the body of the paper.

%%=============================================%%
%% For submissions to Nature Portfolio Journals %%
%% please use the heading ``Extended Data''.   %%
%%=============================================%%

%%=============================================================%%
%% Sample for another appendix section			       %%
%%=============================================================%%

%% \section{Example of another appendix section}\label{secA2}%
%% Appendices may be used for helpful, supporting or essential material that would otherwise 
%% clutter, break up or be distracting to the text. Appendices can consist of sections, figures, 
%% tables and equations etc.

\end{appendices}

%%===========================================================================================%%
%% Bibliography entries                                                                    %%
%%===========================================================================================%%

\begin{thebibliography}{20}
\bibitem{bib1} Husák M, Komárková J, Bou-Harb E, Čeleda P. Survey of attack projection, prediction, and forecasting in cyber security. IEEE Commun Surveys Tuts. 2019;21(1):640--680.

\bibitem{bib2} 2023 Data Breach Investigations Report. Verizon Enterprise Solutions; 2023.

\bibitem{bib3} Whitman ME, Mattord HJ. Principles of Information Security. 6th ed. Boston, MA, USA: Cengage Learning; 2021.

\bibitem{bib4} Barnum S. Standardizing cyber threat intelligence information with the structured threat information expression (STIX). MITRE Corporation, Bedford, MA, USA, Tech. Rep.; 2012.

\bibitem{bib5} Kent K, Chevalier S, Grance T, Dang H. Guide to integrating forensic techniques into incident response. NIST, Gaithersburg, MD, USA, NIST Special Publication 800-86; Aug. 2006.

\bibitem{bib6} Morgan S. Cybercrime to cost the world \$10.5 trillion annually by 2025. Cybercrime Magazine. 2020. Available: https://cybersecurityventures.com/cybercrime-damages-2025/

\bibitem{bib7} Fielding RT. Architectural styles and the design of network-based software architectures. Ph.D. dissertation, Dept. Information and Computer Science, Univ. California, Irvine, CA, USA; 2000.

\bibitem{bib8} Papazoglou MP, Traverso P, Dustdar S, Leymann F. Service-oriented computing: State of the art and research challenges. Computer. 2007;40(11):38--45.

\bibitem{bib9} Framework for Improving Critical Infrastructure Cybersecurity, Version 1.1. NIST, Gaithersburg, MD, USA; Apr. 2018.

\bibitem{bib10} Wagner A, Dulaunoy S, Wagener G, Iklody C. MISP: The design and implementation of a collaborative threat intelligence sharing platform. In: Proc. 2016 ACM Workshop Inf. Sharing Collaborative Security (WISCS), Vienna, Austria; Oct. 2016. p. 49--56.

\bibitem{bib11} Iannacone AL, Bohn M, Nakamura S, Gerth J, Huffer K, Bridges R, Ferragut E, Goodall J. Developing an ontology for cyber security knowledge graphs. In: Proc. 10th Annu. Cyber Inf. Security Res. Conf. (CISR), Oak Ridge, TN, USA; 2015. p. 1--4.

\bibitem{bib12} Khraisat A, Gondal I, Vamplew P, Kamruzzaman J. Survey of intrusion detection systems: Techniques, datasets and challenges. Cybersecurity. 2019;2(1):20.

\bibitem{bib13} Saxe J, Berlin K. eXpose: A character-level convolutional neural network with embeddings for detecting malicious URLs, file paths, and registry keys. arXiv:1702.08568. 2017. Available: https://arxiv.org/abs/1702.08568

\bibitem{bib14} Brown T, Mann B, Ryder N, Subbiah M, Kaplan JD, Dhariwal P, et al. Language models are few-shot learners. In: Proc. Adv. Neural Inf. Process. Syst. vol. 33; 2020. p. 1877--1901.

\bibitem{bib15} Weidinger L, Mellor J, Rauh M, Griffin C, Uesato J, Huang P-S, et al. Ethical and social risks of harm from language models. arXiv:2112.04359. 2021. Available: https://arxiv.org/abs/2112.04359

\bibitem{bib16} Jain R. The Art of Computer Systems Performance Analysis: Techniques for Experimental Design, Measurement, Simulation, and Modeling. New York, NY, USA: Wiley; 1991.

\bibitem{bib17} Challenger J, Dantzig P, Iyengar A. A scalable system for consistently caching dynamic web data. In: Proc. 38th Annu. IEEE/IFIP Int. Conf. Dependable Syst. Netw. (DSN), Seattle, WA, USA; Jun. 1999. p. 1--10.

\bibitem{bib18} Breslau L, Cao P, Fan L, Phillips G, Shenker S. Web caching and Zipf-like distributions: Evidence and implications. In: Proc. IEEE INFOCOM, vol. 1, New York, NY, USA; Mar. 1999. p. 126--134.

\bibitem{bib19} Gamma E, Helm R, Johnson R, Vlissides J. Design Patterns: Elements of Reusable Object-Oriented Software. Reading, MA, USA: Addison-Wesley; 1994.

\bibitem{bib20} Martin RC. Clean Architecture: A Craftsman's Guide to Software Structure and Design. Upper Saddle River, NJ, USA: Prentice Hall; 2017.

\end{thebibliography}

\end{document}
